\chapter{目标试验模拟}

% 本文件内的局部辅助：带垂直偏移的边注，用于错开相邻批注
\providecommand{\mnotev}[2]{\marginnote{\textcolor{black!85}{#2}}[#1]}

如第一部分所述，从观察性数据作因果推断，可以视为一种尝试：模拟一项假想的随机化试验，我们称之为\emph{目标试验}。然而，第一、第二部分只涉及比较时不变处理的、简化了的目标试验。既然我们现在已具备处理时变处理因果推断所需的全部工具，就可以讨论比较\emph{持续处理策略}的现实目标试验了。本章把目标试验的概念推广到持续处理策略，并勾勒一个统一的因果推断框架——无论数据来自随机化实验还是观察性研究。

本章还描述了模拟目标试验时可能关心的因果效应分类，其中包括意向性治疗效应与符合方案效应的观察性对应物。对这些因果效应的有效估计，一般需要时变预后因子与处理的数据，并用 g 方法对这些时变因子作恰当调整。正是 g 方法的发展，才使这里讨论的概念不止于形式操练：只要有数据涵盖所有重要的固定与时变混杂因子，所关心的效应如今就能被有效估计。

\section{意向性治疗效应与符合方案效应}

考虑一项随机化试验：有被某种危险病毒感染风险的个体，被随机分配到“立即接种疫苗，一旦感染即接受实验性抗病毒治疗”的联合处理（$Z=1$），或分配到标准治疗（$Z=0$，其中既不接种疫苗，也不给抗病毒治疗）。图~22.1 用分配处理 $Z$、实际接受的处理 $A$ 与结局（死亡）$Y$ 表示这项试验。对某个个体而言，$Z$ 与 $A$ 的取值可能不同，原因是不依从所分配的处理：有些被分配到疫苗（$Z=1$）的个体可能因拒绝接种而未接受疫苗（$A=0$）；有些被分配到不接种疫苗（$Z=0$）的个体，仍可能在研究之外获得疫苗（$A=1$）。变量 $U$ 表示影响个体接种决定之不可测量的危险因素。

\begin{figure}[htbp]
\centering
\begin{tikzpicture}[causal,every node/.style={transform shape}]
  \node (z) at (0,0) {$Z$};
  \node (a) at (3.0,0) {$A$};
  \node (y) at (6.0,0) {$Y$};
  \node (u) at (1.2,-1.55) {$U$};
  \draw[ed] (z) -- (a);
  \draw[ed] (a) -- (y);
  \draw[ed] (z) to[bend left=40] (y);
  \draw[ed] (u) -- (a);
  \draw[ed] (u) -- (y);
\end{tikzpicture}
\caption{}
\label{fig:22.1}
\end{figure}

如图~22.1 所示，分配处理 $Z$ 可以通过两条不同路径对结局 $Y$ 产生因果效应。第一，处理分配 $Z$ 可能仅仅因为影响实际接受的处理 $A$ 而影响结局 $Y$。被分配到疫苗者更可能接受疫苗，这由 $Z\to A$ 的箭头表示。若接种疫苗对死亡率有因果效应（由 $A\to Y$ 的箭头表示），则分配到疫苗便通过路径 $Z\to A\to Y$ 对结局 $Y$ 有因果效应。

第二，处理分配 $Z$ 也可能通过未经实际接受的处理 $A$ 中介的路径影响结局 $Y$。例如，知道自己被分配了何种处理，可能改变参与者的行为——知道自己被分配了疫苗接种加一种有希望的抗病毒治疗，个体可能对是否被感染不再那么小心。这些行为改变由 $Z$ 到 $Y$ 的直接箭头表示。

\begin{finepoint}[排除限制（再论）]
图~22.1 中箭头 $Z\to Y$ 的存在，表示分配对结局有一项不经过处理的直接效应。当这一箭头存在时，我们说\emph{排除限制}不成立。关于排除限制的正式讨论见技术注记~16.1。

研究者常常试图通过消除箭头 $Z\to Y$ 来部分地“净化”$Z$ 的效应，如图~22.2（与图~16.1 相同）所示，该图刻画了“$Z$ 到 $Y$ 无直接箭头”这一排除限制。为此，他们对参与者及其医生隐瞒所分配处理 $Z$ 的信息。例如，研究者会给被随机分配到 $Z=1$ 者施打疫苗，给被分配到 $Z=0$ 者施打安慰剂（除不含疫苗外完全相同的注射）。由于参与者及其医生不知道所注射的是积极处理还是安慰剂，他们被称为“盲”，该研究被称为双盲安慰剂对照随机化试验。在第 16 章，我们用双盲安慰剂对照随机化试验的概念来引出工具变量的概念。

双盲的处理分配常常不可行。许多研究无法有效设盲，因为没有实际办法施打令人信服的安慰剂（例如在开胸手术中），或者处理的副作用会暴露谁在服用它，等等。此外，当研究者关心的是量化真实世界中的处理效应时——而真实世界中既无设盲也无安慰剂——也不宜设盲（及安慰剂对照）。
\end{finepoint}

\begin{figure}[htbp]
\centering
\begin{tikzpicture}[causal,every node/.style={transform shape}]
  \node (z) at (0,0) {$Z$};
  \node (a) at (3.0,0) {$A$};
  \node (y) at (6.0,0) {$Y$};
  \node (u) at (1.2,-1.55) {$U$};
  \draw[ed] (z) -- (a);
  \draw[ed] (a) -- (y);
  \draw[ed] (u) -- (a);
  \draw[ed] (u) -- (y);
\end{tikzpicture}
\caption{}
\label{fig:22.2}
\end{figure}

因此，分配处理 $Z$ 的因果效应不仅取决于箭头 $A\to Y$ 的强度（实际接受处理的效应），还取决于箭头 $Z\to A$（研究中依从所分配处理的程度）与箭头 $Z\to Y$（同时发生的行为改变）的强度。$Z$ 的效应不是“用 $A$ 治疗的效应”，而是“把参与者分配到用 $A$ 治疗”的效应，或“怀有以 $A$ 治疗的意图”的效应——这正是随机分配 $Z$ 的效应常被称为\emph{意向性治疗效应}的原因。

由于 $Z$ 是随机分配的，分配处理的效应预期没有混杂。在图~22.1 中，$Z$ 到 $Y$ 不存在后门路径，故对分配处理 $Z$，可交换性 $Y^{z}\indep Z$ 预期成立。$Z$ 与 $Y$ 之间的关联意味着 $Z$ 对 $Y$ 有因果效应，无论所有个体是否依从所分配的处理。关联风险比 $\Pr[Y=1\given Z=1]/\Pr[Y=1\given Z=0]$ 等于因果意向性治疗风险比 $\Pr[Y^{z=1}=1]/\Pr[Y^{z=0}=1]$。通过估计 $Z$ 与 $Y$ 之间未经调整的关联来估计意向性治疗效应的分析，称为\emph{意向性治疗分析}。精要~22.2 讨论了几种通常有偏的意向性治疗分析变体。\mnotev{-3.4em}{用统计术语说，意向性治疗分析为“无平均因果效应”这一零假设提供了一项有效——尽管也许检验效能不足——的 $\alpha$ 水平检验，但仅限双盲安慰剂对照的随机化实验。}

现在考虑这样的处理因果效应：若所有个体都依从试验方案所规定的所分配处理，本会观测到的效应，我们称之为\emph{符合方案效应}。本书大部分篇幅都假定对所分配处理的完美依从，从而所有参与者的分配处理 $Z$ 与实际接受处理 $A$ 取值一致。也就是说，我们假定 $U$ 不存在，因而受治疗者（$A=1$）与未受治疗者（$A=0$）可交换，$Y^{a}\indep A$。\mnote{符合方案效应可由对比 $\Pr[Y^{z=1,a=1}=1]$ 与 $\Pr[Y^{z=0,a=0}=1]$ 来定义，或在排除限制下由对比 $\Pr[Y^{a=1}=1]$ 与 $\Pr[Y^{a=0}=1]$ 来定义。正文为记号简便起见采用后者。}

现在考虑这样一种情形：$U$ 表示感染的高风险（1：是，0：否），且 $Z=0$ 组中感染高风险者（$U=1$）倾向于到研究之外寻求接种（$A=1$）。若果真如此，则 $A=1$ 组中高风险个体的比例将高于 $A=0$ 组：$A=1$ 与 $A=0$ 两组不可交换，于是关联风险比 $\Pr[Y=1\given A=1]/\Pr[Y=1\given A=0]$ 不等于（因果）符合方案风险比 $\Pr[Y^{a=1}=1]/\Pr[Y^{a=0}=1]$。如后门路径 $A\leftarrow U\to Y$ 所示，$A$ 对 $Y$ 的效应存在混杂，估计符合方案效应需要调整。也就是说，估计符合方案效应，要求把随机化实验当作一项观察性研究来看待。精要~22.3 描述了几种错过这一点的、量化符合方案效应的传统做法。

\begin{finepoint}[伪意向性治疗分析与改良意向性治疗分析]
意向性治疗分析对意向性治疗效应无偏，因为它纳入了所有被随机化的个体。因此，只纳入被随机化个体中一个子集的意向性治疗分析变体可能是有偏的。

当一些个体未完成随访时，其结局未知，因而分析需要限于随访完整的个体。于是我们只能做\emph{伪意向性治疗分析}
\[
\Pr[Y=1\given Z=1,C=0]/\Pr[Y=1\given Z=0,C=0],
\]
其中 $C=0$ 表示个体在测量 $Y$ 之前一直未被删失。如第 8 章所述，删失可能引入选择偏倚，因而伪意向性治疗估计可能是意向性治疗效应的有偏估计，偏倚方向不定。在存在失访或其他形式删失时，随机化实验的意向性治疗分析需要适当调整选择偏倚。进一步的讨论见 21.5 节与 Little et al. (2012)。

对持续处理策略，一种常见做法是把意向性治疗分析限于至少起始了其所分配策略（例如至少服用了一粒药）的个体。这种做法称为\emph{改良意向性治疗分析}，它只纳入被随机化个体中的一个子集，因而可能对意向性治疗效应有偏。改良意向性治疗分析一般需要调整影响依从性的危险因素。
\end{finepoint}

缺乏混杂，在很大程度上解释了意向性治疗效应在许多随机化实验中为何享有特权地位：“怀有以 $A$ 治疗的意图”的效应也许不是我们想要的效应——“用 $A$ 治疗的效应”，即符合方案效应——但它更容易计算。当一件较不有趣的数量比一件较有趣的数量更容易计算时（这种情况屡见不鲜），我们往往会想方设法为使用那件较不有趣的数量找理由。意向性治疗效应也不例外。下面我们讨论，为什么若干关于意向性治疗效应的著名辩护需要打些折扣。

关于意向性治疗效应的一个常见辩护是它能\emph{保持零效应}。也就是说，若处理 $A$ 对 $Y$ 没有效应，则分配处理 $Z$ 对 $Y$ 也不会有影响。零效应保持是一项关键性质，因为它确保在不存在效应时不会宣称有效应。更正式地说，在严格因果零假设与排除限制之下，可以证明
\[
\Pr[Y=1\given Z=1]/\Pr[Y=1\given Z=0]=\Pr[Y^{a=1}=1]/\Pr[Y^{a=0}=1]=1.
\]
然而，当排除限制不成立时（如图~22.1 所示），这一等式并不成立。在那些情形下——并非双盲安慰剂对照的实验——$A$ 的效应可能为零，而 $Z$ 的效应不为零。要看出这一点，请在脑海中擦去图~22.1 中的箭头 $A\to Y$：仍然存在一条从 $Z$ 到 $Y$ 的箭头。

关于意向性治疗效应的一个相关辩护是：其取值“比符合方案效应的取值更接近零”。其直觉是，若不完美依从导致效应衰减——而非夸大——则意向性治疗风险比 $\Pr[Y=1\given Z=1]/\Pr[Y=1\given Z=0]$ 将落在 1 与符合方案风险比 $\Pr[Y^{a=1}=1]/\Pr[Y^{a=0}=1]$ 的取值之间。于是意向性治疗效应可被解释为符合方案效应的一个下界，即意向性治疗效应是符合方案效应的保守估计。遗憾的是，意向性治疗效应并不总是保守的，因为效应衰减并无保证。见精要~22.4。

\begin{finepoint}[朴素的符合方案分析]
在随机化试验中，试图估计处理 $A$ 的符合方案效应的两种常见做法是“按实际治疗”分析与所谓的“符合方案”分析。

传统的\emph{按实际治疗分析}比较接受治疗者（$A=1$）与未接受治疗者（$A=0$）之间结局 $Y$ 的分布，而不论其处理分配 $Z$。显然，若促使参与者接受治疗的原因与未测量的预后因子 $U$ 相关（如图~22.1 与 22.2），传统的按实际治疗比较就会有混杂。另一方面，考虑这样一种设定：$A$ 与 $Y$ 之间的所有后门路径都可以通过条件于已测量因子 $L$ 而被阻断，如图~22.3 所示。此时按实际治疗分析需要调整因子 $L$。

传统的\emph{符合方案分析}——有时称为\emph{在治疗分析}——只纳入依从研究方案的个体：即 $A=Z$ 的所谓符合方案人群。分析随后仅在这一符合方案人群中比较被分配到治疗者（$Z=1$）与被分配到不治疗者（$Z=0$）之间结局 $Y$ 的分布。也就是说，传统的符合方案分析不过是一项限于符合方案人群的意向性治疗分析。这种限制一般会导致符合方案效应的有偏估计。要看出原因，考虑图~22.4 的因果图，其中包含一个进入符合方案人群的选择指示变量 $S$：若 $A=Z$ 则 $S=1$，否则 $S=0$。除非符合方案分析恰当地测量并调整因子 $L$，否则将出现选择偏倚，因为条件于 $S=1$ 会打开非因果路径 $Z\to A\leftarrow L\leftarrow U\to Y$。

也就是说，按实际治疗分析与符合方案分析都是对随机化实验的观察性分析，且如同任何观察性分析，要获得符合方案效应的有效估计，需要恰当调整混杂与选择偏倚。例子与进一步讨论见 Hernán and Hernández-Díaz (2012)。

\begin{center}
\begin{tikzpicture}[causal,every node/.style={transform shape},scale=0.95]
  \node (z) at (0,0.35) {$Z$};
  \node (l) at (2.3,0.35) {$L$};
  \node (a) at (4.6,0.35) {$A$};
  \node (y) at (6.9,0.35) {$Y$};
  \node (u) at (1.3,-1.1) {$U$};
  \draw[ed] (z) to[bend left=40] (a);
  \draw[ed] (l) -- (a);
  \draw[ed] (a) -- (y);
  \draw[ed] (u) -- (l);
  \draw[ed] (u) -- (y);
\end{tikzpicture}
\captionof{figure}{}
\label{fig:22.3}
\end{center}

\begin{center}
\begin{tikzpicture}[causal,every node/.style={transform shape},scale=0.95]
  \node (z) at (0,0.2) {$Z$};
  \node (l) at (2.3,0.2) {$L$};
  \node (a) at (4.6,0.2) {$A$};
  \node (y) at (6.9,0.2) {$Y$};
  \node (u) at (1.3,-1.2) {$U$};
  \node[draw,inner sep=2.6pt,line width=0.5pt] (s) at (6.0,2.15) {$S$};
  \draw[ed] (z) to[bend left=30] (a);
  \draw[ed] (0.28,0.44) -- (s.west);
  \draw[ed] (a) -- (s.south west);
  \draw[ed] (l) -- (a);
  \draw[ed] (a) -- (y);
  \draw[ed] (u) -- (l);
  \draw[ed] (u) -- (y);
\end{tikzpicture}
\captionof{figure}{}
\label{fig:22.4}
\end{center}
\end{finepoint}

\begin{finepoint}[关于意向性治疗效应的更多误解]
一个常听到的论点是：意向性治疗效应度量了处理在真实世界中的\emph{效果}，因为它纳入了人们不会完美依从所分配处理这一事实。相反，符合方案效应度量的是处理在完美依从下的\emph{效力}。用这套术语，人们常论证“效力”并不反映处理在真实条件下的效应，因而有理由把意向性治疗效应报告为随机化实验的主要发现，因为“效果”才是对处理效应最现实的度量。

这一推理有几处问题。第一，意向性治疗效应度量的是在某一特定实验中所观察到的依从条件下分配处理的效应。真实生活中的实际依从可能不同（例如，研究参与者若受到密切监测，可能依从得更好），而且实际上可能受该特定实验结果的影响（例如，人们得知某治疗有效后会更可能依从）。第二，若目标是效果，我们就应当放弃实施双盲安慰剂对照的随机化临床试验，因为在真实生活中，病人与医生都知道所接受的处理，也不使用安慰剂。真正的效果度量应当纳入源自分配知晓的效应（例如行为改变），而这类效应在双盲随机化实验中被消除。第三，计划依从其医生所开处方的个体，对符合方案效应的兴趣会大于对意向性治疗效应的兴趣。

另一个常见论点是：意向性治疗效应保证保守。这并非在所有设定中都成立。若处理的符合方案效应不单调（即并非对所有个体都同向；见技术注记~5.2），且不依从程度很高，那么符合方案效应可能比意向性治疗效应更接近零。即便对单调效应，在把个体分配到两种积极处理之一的头对头试验中，意向性治疗效应也不一定保守。设想患有疼痛性疾病的个体被随机分配到一种昂贵药物（$Z=1$）或布洛芬（$Z=0$）。目标是确定随 1 年随访后哪种药物导致严重疼痛 $Y$ 的风险更低。研究者并不知道，两种药物减轻疼痛的效力相同，即符合方案风险比为 1。然而，由于一种容易缓解的轻微副作用，对布洛芬的依从恰好低于对昂贵药物的依从。结果，意向性治疗风险比大于 1，研究者便错误地得出结论：布洛芬在减轻严重疼痛方面不如昂贵药物有效。更多细节见 Robins (1998b) 以及 Hernán and Hernández-Díaz (2012) 的讨论。
\end{finepoint}

即便在意向性治疗效应确实保守的情形下，这也许并非好事。设想目标是评价某种处理的安全性：人们可能天真地因为 $Z$ 对不良结局的意向性治疗效应接近零，就断定处理 $A$ 是安全的，即便处理 $A$ 在相当一部分病人中导致了该不良结局。解释或许是：许多被分配到 $Z=1$ 的个体在出现不良结局之前并未服用、或已停止服用该处理。那么，用意向性治疗效应来定义处理的效应就是一种危险的做法。

总之，对于不依从程度较高的随机化试验、以及评价危害而非益处的试验，完全依赖意向性治疗效应估计难以成立。对研究者与决策者（如临床医生、病人）而言，符合方案效应往往是更自然的估计目标。估计符合方案效应，需要在给定已测量协变量的可交换性假设下调整混杂，或在其他假设（例如工具变量估计所需的假设，见第 16 章）下调整。

以上讨论大体围绕时不变处理。当随机化试验研究的是可随时间变化的持续策略（这种情况很常见）时，不依从的概率会大大增加。于是，与符合方案效应相比，意向性治疗效应变得越来越缺乏信息；而符合方案效应被定义为：若每个人在整个随访期间都依从其所分配的处理策略，本会观测到的效应。估计持续策略的符合方案效应——无论在真实的随机化试验中，还是在模拟该试验的观察性分析中——一般都需要 g 方法。

\section{持续处理策略的目标试验}

现在可以讨论比较持续处理策略的目标试验了。由于最终目标是用真实世界的观察性数据来模拟这些试验，我们只考虑特征与真实世界相似的\emph{实用性}试验。具体而言，参与者及其主治医生需要知道所接受的处理（即处理分配不设盲），无人接受安慰剂（即策略 $g$ 与 $g'$ 都只涉及积极处理或不处理），且参与者受监测的频率与强度与试验之外的普通病人相同。当目标是量化现实条件下处理策略的效应时，具有实用性特征的试验更可取。

为把想法固定下来，考虑一项随机化试验，估计抗逆转录病毒治疗对 HIV 感染者 5 年死亡风险的影响。合格参与者——18 岁及以上、无艾滋病、既往未使用过抗逆转录病毒治疗——在随访起点 $k=0$（基线）被随机分配到处理策略 $g$ 或处理策略 $g'$。其随访从分配之时开始，到死亡（所关心的结局）、失访或基线后 60 个月为止，以先发生者为准。

令 $A_k$ 在个体于时点 $k$ 接受治疗时取 1，否则取 0，其中 $k=0,1,2,\ldots,K$，$K=59$。我们的试验把合格个体分配到策略 $g_1$“在随访期间持续接受治疗 $A_k=1$，除非出现禁忌证或毒性”，或策略 $g_0$“在随访期间持续接受治疗 $A_k=0$”。令分配指示变量 $Z$ 在个体被分配到 $g_1$ 时取 1、被分配到 $g_0$ 时取 0。令 $D_k$ 为截至第 $k$ 个月（$k=1,2,\ldots,K+1$）的死亡指示变量（1：是，0：否），$C_k$ 为截至该月的删失指示变量（1：是，0：否）。

现在我们来定义具有持续处理策略的随机化试验中的意向性治疗效应与符合方案效应。持续策略的其他对比——称为\emph{直接效应}——见技术注记~22.1。\mnote{前几章考虑的是处理对随访结束时测量的结局 $Y$ 的因果效应。在这项试验中，结局是失效时间，即死亡时间（见技术注记~21.10）。}

意向性治疗效应是如下静态策略的对比
\begin{itemize}
\item $(z=1,\bar c_K=\bar 0)$：在基线被分配到策略 $g_1$，并留在研究中直至随访结束；
\item $(z=0,\bar c_K=\bar 0)$：在基线被分配到策略 $g_0$，并留在研究中直至随访结束。
\end{itemize}
于是，时点 $k$ 的意向性治疗效应可以表示为死亡的反事实风险之对比
\[
\Pr\bigl[D_k^{z=1,\bar c_k=\bar 0}=1\bigr]-\Pr\bigl[D_k^{z=0,\bar c_k=\bar 0}=1\bigr],
\]
即：在分配到策略 $g_1$ 与 $g_0$ 之下、若截至时点 $k$ 无人失访（$\bar c_k=\bar 0$）时的对比。

在某些随机化试验中，对处理策略的分配与治疗的起始同时发生。也就是说，所有被分配到策略 $g_1$ 的个体都在时点 0 开始接受治疗，无论其在基线之后是否继续服用；而所有被分配到策略 $g_0$ 的个体在时点 0 都不接受治疗，无论其在基线之后是否开始服用。在这些情形下，意向性治疗效应不仅是分配的效应，也是治疗起始的效应
\[
\Pr\bigl[D_k^{a_0=1,\bar c_k=\bar 0}=1\bigr]-\Pr\bigl[D_k^{a_0=0,\bar c_k=\bar 0}=1\bigr].
\]

如同任何随机化试验，一些参与者会违背方案，不依从所分配的策略。在随访期间，一些被分配到 $g_1$ 的个体会无临床理由地停止治疗，一些被分配到 $g_0$ 的个体会开始治疗，一些个体会使用未经批准的伴随治疗，等等。意向性治疗效应对这些方案违背不作区分，而它们都是基线之后所作决定的结果。这种不区分意味着，意向性治疗效应的大小可能严重依赖各项试验实施过程中出现的特定方案违背模式。两项方案相同但在不同环境中实施的试验，可能得到不同的意向性治疗效应估计，而二者都没有偏倚。鉴于意向性治疗效应的局限，我们希望用符合方案效应来补充它。

符合方案效应由如下干预下结局分布的对比来定义：
\begin{itemize}
\item 在基线 $k=0$ 与随访结束之间持续接受处理策略 $g_1$；
\item 在基线 $k=0$ 与随访结束之间持续接受处理策略 $g_0$。
\end{itemize}
于是，时点 $k$ 的符合方案效应可以表示为死亡的反事实风险之对比
\[
\Pr\bigl[D_k^{g_1,\bar c_k=\bar 0}=1\bigr]-\Pr\bigl[D_k^{g_0,\bar c_k=\bar 0}=1\bigr],
\]
即：若截至时点 $k$ 无人失访（$\bar c_k=\bar 0$）时依从策略 $g_1$ 与依从策略 $g_0$ 的对比。

\begin{techpoint}[受控直接效应]
考虑把中介 $M$ 设定为某一特定取值时，处理 $A$ 对结局 $Y$ 的平均因果效应。我们把这一数量称为 $A$ 对 $Y$ 不经过 $M$ 的直接效应。若中介 $M$ 可取两个值（0 或 1），我们就可以定义 $M$ 设定为 1 时 $A$ 对 $Y$ 的直接效应，以及 $M$ 设定为 0 时 $A$ 对 $Y$ 的直接效应。在加法尺度上，这两个直接效应分别由反事实差值
\[
\E[Y^{a=1,m=1}]-\E[Y^{a=0,m=1}]
\quad\text{与}\quad
\E[Y^{a=1,m=0}]-\E[Y^{a=0,m=0}]
\]
定义。这两个直接效应常被称为\emph{平均受控直接效应}，原则上可以通过对处理 $A$ 与中介 $M$ 都作序贯随机化的实验来识别，或用观察性数据模拟这样的目标实验来识别。技术注记~22.2 描述了另一些并不存在相应目标实验的直接效应。

设想我们实施一项随机化实验：参与者在基线被随机分配到处理 $A=1$ 或 $A=0$，并在基线后一个月被随机分配到 $M=1$ 或 $M=0$。于是所有个体都将落入四组之一：$(A=1,M=1)$、$(A=1,M=0)$、$(A=0,M=1)$ 或 $(A=0,M=0)$。所关心的结局 $Y$ 在第 3 个月对所有人测量（为简单起见，假定无人失访或死亡）。这一研究设计使我们能一致地估计受控直接效应，因为对 $A$ 与 $M$ 都作随机化，确保了反事实数量 $\E[Y^{a,m}]=\Pr[Y^{a,m}=1]$ 可由观测风险 $\Pr[Y=1\given A=a,M=m]$ 一致地估计。

只要一致性、正值性与可交换性这三个可识别条件对 $A$ 与 $M$ 都成立，受控直接效应也可以在观察性研究中被有效估计。对这些可识别条件的精确刻画其实已在第 19 章给出，因为受控直接效应不过是持续于时间的处理策略之对比的一个特例。要看出这一点，只需把上式中的 $A$ 与 $M$ 换成 $A_0$ 与 $A_1$。更一般地说，处理 $A$ 与中介 $M$ 本身都可以是时变的。
\end{techpoint}

明智的试验方案不会规定无论个体发生什么都要继续治疗。例如，我们的持续治疗策略 $g_1$ 规定：一旦出现禁忌证或毒性，即须停止治疗。也就是说，符合方案效应一般涉及\emph{动态}策略（“这样做；若发生 X，则改做那样”）之间的比较，而非\emph{静态}策略（“无论发生什么都这样做”）之间的比较。

有时研究方案并未明确说明处理策略的动态性质。例如，方案可能把策略 $g_1$ 简化为“在随访期间持续接受治疗 $A_k=1$”，而不明确写出“一旦出现禁忌证或毒性”即须停药。这种对策略 $g_1$ 的简化描述可能引起误解。具体而言，一个被分配到 $g_1$、因毒性而停药的人，不应被标签为不依从策略 $g_1$ 的人。事实上，此人完全依从了（本应）在方案中写明的策略 $g_1$。当在真实世界中别无选择时，停止原分配的治疗或开始其他有医疗指征的治疗，不可能被视为偏离方案。由于符合方案效应由现实策略的对比来定义，它对于旨在为真实世界决策提供证据的因果推断研究尤为相关。\mnote{理想情况下，为避免在何为、何不为随访期间的不依从上产生混乱，方案应完整规定所关心的处理策略。这样符合方案效应才有良好定义 (Hernán and Robins, 2017)。}

事实上，符合方案效应常常是推断的隐含目标。例如，研究者常质疑研究中实施的干预相对于方案所述干预的忠实程度，并称存在“偏倚”。这种说法表明，研究者真正关心的是比较随访期间按方案实施的干预（即符合方案效应），而不是基线上把个体分配到干预的效应（即意向性治疗效应），因为基线之后的不依从不可能使基线分配的效应产生偏倚。

\begin{techpoint}[纯直接效应与主层直接效应]
除技术注记~22.1 所述的受控直接效应之外，把潜在中介 $M$ 设定为某一特定取值时，处理 $A$ 对结局 $Y$ 的平均直接效应还有其他定义。

$A$ 对 $Y$ 不经过 $M$ 的\emph{纯直接效应}（也称\emph{自然直接效应}），是把 $M$ 的取值设定为“若 $A$ 被设定为 0 时 $M$ 本会取的取值”时 $A$ 对 $Y$ 的平均因果效应，即把 $M$ 设定为 $M^{a=0}$（对某些个体为 1，对另一些为 0）。纯直接效应由对比
\[
\E[Y^{a=1,M^{a=0}}]-\E[Y^{a=0,M^{a=0}}]
\]
定义，它是一个\emph{跨世界}数量，因为 $\E[Y^{a=1,M^{a=0}}]$ 包含一个同时以 $a=1$ 与 $a=0$ 为下标的反事实结局。因此，纯直接效应既无法从对 $A$、$M$ 或二者的随机化实验中识别，也无法在 FFRCISTG 模型下从观察性数据中识别（见技术注记~6.2）。尽管如此，估计纯直接效应常常是因果中介分析的目标，因为总处理效应可以分解为纯直接效应与总间接效应。纯直接效应由 Robins and Greenland (1992) 引入；Pearl (2001) 将其改名为自然直接效应，并证明对某些因果图，纯直接效应可以在其 NPSEM-IE 模型下从观测数据中识别，因为与 FFRCISTG 模型不同，NPSEM-IE 模型假定了无法检验的跨世界独立性，而这些独立性无法被对 $A$、$M$ 或二者的随机化实验所驳倒。综述见 VanderWeele (2015) 的著作。

$A$ 对 $Y$ 在 $M$ 的取值被设定为 $m$ 时的\emph{主层直接效应}，是在“无论 $A$ 取何值、$M$ 的取值都会等于 $m$”的那部分总体子集中，$A$ 对 $Y$ 的平均因果效应，即在 $M^{a=0}=M^{a=1}=m$ 的子总体中的平均因果效应。于是主层直接效应定义为
\[
\E\bigl[Y^{a=1,m}\given M^{a=0}=M^{a=1}=m\bigr]-\E\bigl[Y^{a=0,m}\given M^{a=0}=M^{a=1}=m\bigr].
\]
有意思的是，它等于
\[
\E\bigl[Y^{a=1}\given M^{a=0}=M^{a=1}=m\bigr]-\E\bigl[Y^{a=0}\given M^{a=0}=M^{a=1}=m\bigr].
\]
因此，主层直接效应不涉及联合反事实 $Y^{a,m}$，只涉及总体某一子集中的反事实 $Y^{a}$；在这个意义上，它们是该子总体中处理的\emph{总}效应（而非直接效应）。由此可见，与受控直接效应或纯直接效应不同，主层直接效应不要求对 $M$ 的干预有良好定义。当 $A$ 对几乎所有个体都影响 $M$ 时，主层直接效应几乎没有政策相关性，因为它只适用于 $M^{a=0}=M^{a=1}$ 的那一小部分总体。实践中常把 $M$ 粗化（通常粗化为二值指示变量）以增大主层的规模，但粗化本身可能使主层直接效应的科学相关性降低 (Robins et al. 2007)。主层直接效应由 Robins (1986) 引入，并由 Rubin (2004) 推广。Frangakis and Rubin (2002) 循 Robins (1986) 的思路，把主层这一概念用作处理竞争事件的工具。在第 23 章，我们考虑一种干预主义的中介理论 (Robins and Richardson 2010)，它提供了又一种直接效应。
\end{techpoint}

最后，考虑接受研究方案所规定之外的干预的效应。假设在我们的试验进行期间，一种共识开始形成：策略 $g_0$“在随访期间持续接受治疗 $A_k=0$”劣于策略 $g_1$。于是，一些医生开始在临床病程恶化、CD4 细胞计数（$L_k$）首次降至 200 个/$\mu$L 以下时建议开始治疗。结果，试验中许多被分配到策略 $g_0$ 的个体实际上遵循了修改后的策略 $g_0'$“在随访期间持续接受治疗 $A_k=0$，但在 $L_k<200$ 之后转为治疗 $A_k=1$”。如下干预下结局分布的对比
\begin{itemize}
\item 在基线 $k=0$ 与随访结束之间持续接受处理策略 $g_1$；
\item 在基线 $k=0$ 与随访结束之间持续接受处理策略 $g_0'$，
\end{itemize}
既不对应意向性治疗效应，也不对应原来的符合方案效应。相反，它问的是这样一个假想目标试验中的符合方案效应：个体被随机分配到策略 $g_1$ 或 $g_0'$。

这个例子说明，与原来的符合方案效应不相符的、所关心的因果效应，如何可以被设想为某些目标试验的符合方案效应，而这些目标试验可以用该随机化试验的数据来模拟。有意思的是，若所关心的策略与实际试验中的策略不同，那么让实际试验的所有参与者都依从方案所规定的策略，反而是不利的。完全依从意味着试验数据无法用来模拟方案不同的目标试验（因为实际数据中无人遵循新目标试验的方案）。例如，一项完全依从的随机化试验把 HIV 感染者分配到不同的 CD4 细胞计数阈值开始抗逆转录病毒治疗，这项试验对于模拟“把个体分配到持续治疗或不治疗”的试验几乎无用，反之亦然。正是这种不依从，才使我们能够用某项随机化试验的数据去模拟其他随机化试验，回答不同的、或许更相关的因果问题。

在具有持续处理策略的随机化试验中，估计符合方案效应所引出的问题，与观察性研究中任何持续策略的比较相同。正如我们后面所讨论的，符合方案效应的有效估计一般要求试验研究者收集随机化之后关于策略依从性、以及与依从性相关的时变预后因子的数据。\mnote{关于随机化试验中符合方案效应估计的更多细节，见 Hernán and Robins (2017)。}

\section{模拟采用持续策略的目标试验}

若无法实施实用性随机化试验，我们可以尝试通过分析现有观察性数据来模拟它。此时我们把该试验称为观察性分析的\emph{目标试验}。

规定目标试验的方案，是澄清我们希望观察性分析回答哪个所关心因果问题的有用手段。至少，我们需要规定方案的以下关键组成部分：合格标准、随访的起点与终点、处理策略、所关心的结局、因果对比与数据分析计划。注意，精确规定目标试验的方案可能需要对现有数据作一些探查。例如，只有在确定数据包含 HIV 诊断信息之后，我们才能合理地提议模拟一项针对 HIV 感染者的目标试验。

上一节所述随机化试验因果效应的对应物，可以为模拟目标试验的观察性分析提出。在现有数据的观察性分析中，模拟意向性治疗效应很少可行，因为对处理策略的实际分配未知。在我们的例子中，意向性治疗效应最接近的观察性对应物，是对不同处理策略之\emph{起始}的比较。对起始者的比较，相当于目标试验中分配与治疗起始总在基线同时发生（无论个体在基线之后是否继续该策略）时的意向性治疗分析。我们可以用如下假想干预下结局分布的对比，来定义意向性治疗效应的这一观察性对应物
\begin{itemize}
\item 在基线起始治疗 $A_0=1$，并留在研究中直至随访结束；
\item 在基线起始治疗 $A_0=0$，并留在研究中直至随访结束。
\end{itemize}
于是，时点 $k$ 的意向性治疗效应之观察性对应物可以表示为反事实风险之对比
\[
\Pr\bigl[D_k^{a_0=1,\bar c_k=\bar 0}=1\bigr]-\Pr\bigl[D_k^{a_0=0,\bar c_k=\bar 0}=1\bigr].
\]
与按所分配策略定义分组的真正意向性治疗效应不同，这一对比按各策略的起始来定义分组。若我们在随机化试验中使用这一对比，就会把所有在基线未服用任何剂量治疗的人归入同一组，无论他们被分配到策略 $g_1$ 还是 $g_0$。若治疗的起始发生在分配到治疗之后不久，我们的观察性对应物大致保留了意向性治疗效应的一项关键特征：对比由基线之后不久发生的干预来定义。\mnote{若我们有抗逆转录病毒治疗的处方（而非发药）数据，那么按基线时是否获得治疗处方来比较各组，会在某种程度上更接近目标试验中的意向性治疗分析。}

另一方面，符合方案效应的观察性对应物，其定义与目标试验中的完全相同。在随机化试验中，我们区分了原来的符合方案效应与其他目标试验中的符合方案效应。在观察性研究中，这种区分没有必要，因为若无预先规定的方案，每个符合方案效应都对应于某个特定的目标试验。一般说来，我们只能用观察性数据模拟那些其预期干预确实被研究中至少一部分个体遵循的目标试验。不过在有些情形下，研究者可能愿意使用建模（例如剂量—反应结构模型）来外推到数据中实际存在的干预之外。\mnote{比较起始者的观察性分析，等价于精要~22.2 所述的改良意向性治疗分析。}

用目标试验中的因果效应来定义观察性研究中的因果效应，迫使我们明确说明所比较的策略。这种对处理策略的明确规定防止了偏倚，因为它使某些数据分析显然包含着无法转译为假想干预之对比的比较。因此，当分析目标是帮助决策者在若干行动方案中作出选择时，这类数据分析应当避免，如我们在 3.5 节与 3.6 节所讨论的。

明确界定观察性分析中处理策略的另一个好处是清晰。如精要~22.4 所述，一些研究者坚持把因果效应分为“效力”（宽泛地定义：在完美条件下本会观测到的处理效应）或“效果”（宽泛地定义：在现实条件下本会观测到的处理效应）。有时随机化试验中的意向性治疗效应被解释为处理的效果，而同一试验中的符合方案效应被解释为处理的效力。另一些时候，随机化试验中的意向性治疗效应被解释为效力（即便在不完美条件下，如不依从），而模拟该试验的观察性研究中的符合方案效应被解释为效果（即便在完美依从下）。也就是说，尤其在涉及长期持续策略的情形中，“效果”与“效力”这两个标签含混不清：往往难以论证说，不依从情形下的意向性治疗效应，或真实世界情形下的符合方案效应，度量的是完美条件下处理的因果效应。

与其固守人为的效力—效果二分，不如承认：一切因果效应都位于效果连续统的某处。此时，明确定义界定所关心因果效应的处理策略更有信息量，因为决策者需要的是关于定义良好的干预之效应的信息。

\section{时点零}

目标试验模拟的一个关键组成部分，是在观察性分析中确定随访的起点，也称基线或\emph{时点零}。合格标准必须在该时点满足，而不能在此之后满足；研究结局从该时点之后开始计数，而不能在此之前计数。

在随机化实验中，每个个体的时点零，是其满足合格标准并被分配到某一处理策略的时刻。例如，在我们关于抗逆转录病毒治疗的随机化试验中，时点零是分配处理策略的时刻（随机化时刻），它通常略早于、或同时于治疗起始。我们不会把随访起点定在比如处理分配之前或之后的 2 年。在随机化之前开始不合情理，因为处理策略尚未分配，合格标准也尚未定义，更谈不上满足；在随机化之后开始则可能有偏，因为试验头两年内的死亡会被排除在分析之外，处理的任何短期效应都会被漏掉。更有问题的是，若处理确实有短期效应，那么在分配到积极治疗的那一组中，更多易感个体到第 2 年时已经死亡，而另一组则不然。两年后易感个体比例的这种差异，摧毁了随机化所达到的基线可比性，为选择偏倚打开了大门。

关于时点零的同样规则适用于观察性分析与随机化试验，理由也相同。一般说来，观察性分析中的随访应当在目标试验中本应开始随访的时点开始。否则，效应估计可能难以解释，并因受处理影响的选择而有偏。尽管如此，在用于因果推断的观察性研究中，对目标试验时点零的模拟错误非常普遍。这些错误源于两个常见问题：1）有时时点零的选择不唯一；2）有时处理策略无法在时点零被唯一分配。下面我们分别描述这两个问题的解决办法。\mnote{例：关于绝经后激素治疗与心脏病，观察性研究与随机化试验的估计值之间那场广为人知的差异，部分原因在于前者对时点零处理不当 (Hernán et al. 2008)。}

第一，时点零不唯一的问题。按照个体一生中满足合格标准的次数，考虑两种情形：

\begin{enumerate}
\item \emph{合格标准只能在单一时刻满足。}这是最简单的设定。随访在该唯一满足合格标准的时刻开始。例如，考虑一项在 HIV 感染者中的研究，比较 CD4 细胞计数首次降至 500 个/$\mu$L 以下时立即开始抗逆转录病毒治疗，与 CD4 细胞计数首次降至 350 个/$\mu$L 以下时延迟开始。合格个体的随访从其 CD4 细胞计数首次降至 500 以下的时刻开始。
\item \emph{合格标准可在多个时刻满足。}这种情形常引起混乱。例如，考虑一项研究，比较无慢性病史、且过去两年未使用激素治疗的绝经后女性中，开始与不开始激素治疗。若一位女性在 51 岁至 65 岁之间持续满足这些合格标准，她的随访应当何时开始？51 岁、52 岁、53 岁……？在目标试验中，一位女性在其一生中的多个时刻都有资格被招募，即她有多个合格时刻。
\end{enumerate}

在有多个合格时刻的设定中，从个体的各个合格时刻里选取时点零，有几种备选做法。可以把时点零选为：a）第一个合格时刻；b）随机选取的一个合格时刻；c）每一个合格时刻，等等。策略 c）要求模拟多个\emph{相继}目标试验，每个试验的随访起点不同。相继试验的数目取决于收集处理与协变量数据的频率：

\begin{itemize}
\item 若按预先规定的时点固定安排数据收集（例如像许多流行病学队列那样每两年一次），则在每个预先规定的时点开始模拟一项新试验。
\item 若数据收集按个体特定的时间安排（例如电子病历），则选定一个固定的时间单位（如一天、一周或一个月），并在每个时间单位开始模拟一项新试验。
\end{itemize}

从统计角度看，相继模拟策略 c）可能比前述策略更有效，因为它利用了更多可得数据。然而，由于个体可能被纳入多个目标试验，需要对效应估计的方差作恰当调整。这可以通过对整个分析做自助法来实现。\mnote{时间单位的不偏选择可因研究而异。例如，考虑这样一项研究：对许多个体而言，时变处理与混杂因子在一周内变化不止一次。此时以一周或一个月为时间单位将引入偏倚。若以天为时间单位，这一偏倚可以消除。若没有处理与混杂因子的逐日数据，偏倚就无法完全校正。}

第二，我们来谈如何应对无法给每个个体分配唯一处理策略的问题。考虑这样一项目标试验：CD4 细胞计数刚降至 500 个/$\mu$L 以下的个体被分配到以下策略之一：（1）立即开始治疗；（2）CD4 细胞计数降至 350 个/$\mu$L 以下时开始治疗；（3）CD4 细胞计数降至 200 个/$\mu$L 以下时开始治疗。用观察性数据模拟这项目标试验时，我们会发现有些个体在时点零（即其 CD4 细胞计数首次降至 500 个/$\mu$L 以下的时刻）就开始了治疗，因而我们会把他们分配到策略（1）。然而，另一些个体在时点零并未开始治疗，这意味着他们的数据在基线上同时与遵循策略（2）和策略（3）相容。我们应当把他们分配到哪个策略？

一种可能是随机选定单一策略并把他们分配到该策略，但这样在统计上无效率。另一种可能是，在数据中为这些个体各造两个完全相同的副本——\emph{克隆}——并把两个克隆分别分配到不同策略。克隆在其数据不再与所分配组别相容的时刻被删失。例如，若该个体在 CD4 降至 350 时未开始治疗，则被分配到“CD4 细胞计数降至 350 时开始治疗”的克隆在该时刻被删失。这种很可能具有信息性的删失所引入的潜在偏倚，需要通过 IP 加权调整时变因子来校正。重要的是，若该个体在两个克隆被删失之前已经死亡，则两个克隆都会死亡，因而这一死亡会被同时分配给两个策略。事件的这种双重分配，防止了如下偏倚：若等待期内发生的事件被系统性地只分配给两个策略之一，就会产生这种偏倚。

\begin{finepoint}[宽限期]
考虑一项观察性研究，比较在时点零立即开始抗逆转录病毒治疗与延迟开始。在真实世界中，抗逆转录病毒治疗不可能在分配的同一天恰好开始。视医疗体系而定，完成必要的临床与行政程序、并使病人得到充分告知，可能需要数周或数月。因此，研究者需要规定一个\emph{宽限期}（比如 3 个月），在此期间开始治疗仍被视为“立即”。否则，研究将是在估计那些在现实中不常发生、或无法在实践中成功实施的策略的效应。

使用宽限期的一个后果是：在宽限期内，一个个体的观测数据与不止一个策略相容。例如，在上述研究中，引入 3 个月宽限期意味着干预被重新定义为“在时点零之后 3 个月内开始治疗”对“从不开始治疗”。因此，在基线后第 3 个月开始治疗者，其数据在第 1、2 个月与两个策略都相容。若其中一些人在那两个月内死亡，我们应当把这些死亡分配给哪个策略？如正文所述，我们可以把这些个体随机分配到两个策略之一，或者更好的是，为每个个体造两个克隆，并把两个克隆分别分配到不同策略。当克隆的数据不再与其所分配策略相容时，克隆即被删失。例如，若该个体在第 3 个月开始治疗，则被分配到“3 个月之后开始”的克隆将在该时点被删失。删失引入的潜在偏倚可以通过 IP 加权处理。

当把宽限期与克隆、删失一起使用时，意向性治疗效应无法估计，因为几乎每个人都会为每个处理策略贡献一个克隆。由于每个个体在基线上都被分配到所有策略，基于基线分配的对比（即“意向性治疗分析”）将比较结局实质相同的组。因此，基线设宽限期的分析着眼于估计某种形式的符合方案效应，因而一般需要纳入恰当的调整。

最后注意，有良好定义的、带宽限期的起始策略应规定在宽限期内开始治疗的时机。细节见 Cain et al. (2010) 的附录。
\end{finepoint}

再强调一次，由于个体可能通过其克隆在分析中被多次纳入，需要通过自助法对效应估计的方差作恰当调整。当合格标准可在多个时刻满足时，克隆 + 删失 + 加权程序可以与相继目标试验模拟结合使用。精要~22.5 描述了如何处理那些可在时点零之后的宽限期内（而非恰在时点零）起始的策略。\mnote{关于克隆 + 删失 + 加权程序的描述，见 Robins et al. (2008) 与 Cain et al. (2010)。相关工作见 van der Laan and Petersen (2007)。}

\section{用数据回答“如果”问题的统一进路}

本书描述并整合了两个因果推断框架：反事实与因果图。明确的目标试验模拟把两个框架都复述了一遍，并使它们落脚于可操作的因果推断。通过围绕一个极为熟悉的科学概念——实验——来组织因果推断，目标试验框架帮助研究者运用其领域知识，把定义良好的因果推断问题表述清楚。一旦因果问题被陈述得几乎没有歧义，研究设计与数据分析便自然展开。

目标试验框架适用于许多学科中广泛的因果问题，无论各领域偏好何种术语与方法。例如，经济学家常把混杂与条件可交换性分别称为遗漏变量偏倚与可观测变量上的选择，而传统社会科学家不大使用 g 方法，因为他们的因果问题通常不是围绕时变处理组织的。但与所有关心因果推断的健康与社会科学家所面对的根本任务相比，这些学科差异是表面的：他们都需要把因果问题表述为定义良好的反事实之对比。目标试验框架通过帮助界定导致定义良好反事实的定义良好的干预，促进了这一任务。

目标试验框架还提供了一种共同语言，用以统一随机化试验与观察性研究的因果分析。除基线随机化之外，模拟目标试验的观察性数据分析与真正的随机化试验分析之间没有其他必要的差别（见精要~22.6）。也就是说，随机化试验可以被视为一项有基线随机化的随访研究，而观察性纵向数据则是一项无基线随机化的随访研究。

随着越来越多的随机化实验试图在真实世界环境中估计长期持续处理策略的效应，有基线随机化与无基线随机化的随访研究之间的相似性在健康与社会科学中日益明显。这些研究与 20 世纪中叶把随机化试验推上研究设计层级顶端的那种严格受控环境中的短期实验相去甚远。对于涉及长期持续处理策略的因果问题，可能大幅偏离方案的随机化实验（例如对所分配策略的不完美依从、失访）会受到混杂与选择偏倚的影响，而我们过去习惯于把这类偏倚只与观察性研究联系在一起。

特别是，在估计符合方案效应时，随机化试验与观察性研究都可能需要调整那些预测退出（选择偏倚）与治疗（混杂）的时变预后因子。也就是说，本书所述的因果推断方法同样适用于随机化试验与观察性研究的符合方案分析。而且，正如从观察性数据估计因果效应并不保证成功一样，随机化试验的符合方案效应估计也可能有偏。\mnote{观察性研究中的时变混杂，是一种与随机化试验中非随机不依从结构相同的偏倚。}

鉴于这些相似性，人们或许会预期随机化实验与观察性研究应以相似方式分析，只是在观察性分析中为估计意向性治疗效应的对应物而调整基线混杂因子。然而在实践中，随机化实验与观察性研究的典型分析截然不同，这既令人困惑，也如我们下面所论证的，成问题。

一个自然的问题是：随机化试验中常用的“意向性治疗分析”与所谓“符合方案分析”，是否分别有效地估计了意向性治疗效应与符合方案效应。典型的意向性治疗分析比较随机化各组之间结局的分布，而不对混杂或选择偏倚作任何形式的调整。不对基线混杂作调整，其正当性来自随机化：随机化各组预期可交换，因为若两组都被分配到同一处理策略，预期它们有相同的结局风险。也不需要对随机化之后的混杂（例如由不依从引起）作调整，因为同样地，基线分配的效应不可能存在随机化之后的混杂。

\begin{finepoint}[随机化实验与观察性研究的数据有何不同？]
只有三件事区分随机化实验与观察性研究的数据。在随机化实验中，(i) 由于随机化，预期没有基线混杂；(ii) 随机化概率已知；(iii) 每个个体在基线上被分配到哪一处理策略已知。

观察性分析可以模拟 (i)，条件是测量并恰当调整一组充分的协变量；也可以模拟 (ii)，条件是给定过去的治疗分配模型被正确设定。有意思的是，(iii) 对估计符合方案效应并非必需，无论在随机化实验还是观察性研究中，因为有效估计量（它们是充分统计量的函数）并不使用这一信息。也就是说，分析者不需要知道所比较的是哪些策略，更不需要知道谁被分配到哪个策略：在随机化试验中，你可以把随机化分配从数据集中删去，只要测量了一组充分的混杂因子，仍能估计符合方案效应。在一项依从完美的动态策略试验中，充分集就是策略在分配治疗时所使用的全部时不变与时变协变量 (Robins 1986)。
\end{finepoint}

然而，基线随机化无法确保随机化之后失访者与未失访者之间的可交换性。由于界定意向性治疗效应的策略要求个体留在研究中直至其结局变量可被确定，在未失访者中计算的意向性治疗效应估计，若预后因子影响、或与失访的差异有关，就可能受随机化之后选择偏倚的影响。因此，意向性治疗效应的有效估计，可能需要一项对随机化之后（时变）预后因子作调整的“意向性治疗分析”，以消除失访带来的选择偏倚。例如，在一项针对 HIV 病人的抗逆转录病毒治疗随机化试验中，若退出研究的概率受症状出现或其他结局危险因素的影响，就需要 g 方法。\mnote{精要~22.2 把连调整选择偏倚都不尝试的意向性治疗分析称为伪意向性治疗分析。}

除主要的意向性治疗分析外，许多随机化试验还报告所谓符合方案分析的结果，该分析仅限于依从研究方案所规定指示的个体，如精要~22.3 就时点干预所述。对持续处理策略，个体在首次偏离方案时被删失。也就是说，每个时点上剩余的符合方案总体，是仍依从方案的那组个体。不作任何形式的调整。这种未调整的分析有三点可疑。\mnote{精要~22.3 把连调整混杂都不尝试的符合方案分析称为朴素的符合方案分析。}

第一，与意向性治疗分析一样，可能因失访的差异而存在选择偏倚。若如此，就需要通过 g 方法调整基线之后（时变）的危险因素。

第二，该分析部分地无视了随机化分组，因此在一种策略下仍留在方案内的个体子集，可能与另一策略下仍在方案内的子集不可交换。也就是说，这种“符合方案分析”类似于观察性分析，因而需要 g 方法来调整由影响“是否继续留在方案内”这一决定的时变危险因素所致的偏倚。工具变量估计（第 16 章）有时可用于有效估计符合方案效应，而不显式调整任何变量，但这些方法的有效性取决于拥有有效的工具以及较强的建模假设。某些形式的工具变量估计是 g 估计的特例（见技术注记~16.6）。

第三，这种传统的符合方案分析忽略了所比较的持续处理策略是动态策略。一个常见错误是把停止治疗的个体当作删失，仿佛停止治疗按定义就是偏离方案——这也是该分析又称\emph{在治疗分析}的原因。上文已讨论过，因毒性或禁忌证而停药的人并未偏离方案，因此不应被删失。\mnote{对失效时间结局，当处理对结局有因果效应时，总是需要 g 方法。原因是处理 $A_k$ 通过其对时变指示变量 $D_{k+1}$ 的效应，影响时点 $k$ 之后的所有变量，如技术注记~21.10 所讨论。}

以上全部考量都适用于随机化试验与旨在模拟目标试验的观察性数据的分析。当目标是估计符合方案效应或其观察性对应物时，随机化试验与观察性研究的分析应当完全相同。如果我们觉得在观察性研究的分析中有必要调整时变混杂与选择偏倚，那么我们也应同样觉得在随机化试验的分析中有必要调整随机化之后的混杂与选择偏倚。用 g 方法调整时变因子，对于随机化试验与观察性研究的符合方案分析一般都将必不可少。目标试验框架与 g 方法使对持续处理策略实施统一的因果推断成为可能。历史上，就作出因果推断与辅助决策而言，随机化实验一直被认为远优于观察性研究。遗憾的是，随机化实验并非总是可行，因为它们可能昂贵、不可行、不合伦理，或只是来不及支持一项紧迫的决策。因此，尽管我们十分看重随机化的好处，但许多决策确实需要在缺乏随机化试验证据的情况下作出。当我们无法实施本可回答因果问题的随机化实验时，就退而尝试用观察性数据模拟它。因此，采用稳妥的方法来设计并分析观察性研究十分重要。把目标试验明确说出来，是朝这个方向迈出的一步。当目标是辅助决策时，对现有观察性数据的分析需要明确地模拟一项试验，并就该分析对其目标试验的模拟程度接受评价。\mnote{在极少数情况下，基于高质量的随机化试验所作的决策，可能反而不如基于严重混杂的观察性数据所作的决策，如精要~22.7 与 22.8 所述。}

\begin{finepoint}[随机化试验与观察性研究的一个反直觉比较]
一种未经检验的非处方处理 $A$ 被某国许多肺癌患者使用。这令该国药品监管机构担忧，于是它出资，对接下来一年内诊断为肺癌者中的一个 20\% 随机样本，开展了一项关于 $A$ 的双盲安慰剂对照随机化试验。所有试验参与者都依从了所分配的处理。60 个月死亡率在治疗组为 55\%，在安慰剂组为 45\%，如下表所示：

\begin{center}
\begin{tabular}{lcc}
\toprule
& $A=0$ & $A=1$\\
\midrule
$Y=1$ & 450 & 550\\
$Y=0$ & 550 & 450\\
\midrule
& 1000 & 1000\\
\bottomrule
\end{tabular}
\end{center}

结果，监管机构禁用了该处理。后来，有人对未选入试验的那 80\% 肺癌患者做了一项观察性研究。该研究发现，在与试验相同的时期内，治疗者与未治疗者的死亡率均为 0\%。观察性数据与随机化试验数据如何才能调和？

首先让我们记住，个体可按反事实类型分类：1）“注定者”（$Y^{a=0}=Y^{a=1}=1$）；2）“受害”者（$Y^{a=0}=0$，$Y^{a=1}=1$）；3）“受益”者（$Y^{a=0}=1$，$Y^{a=1}=0$）；4）“免疫”者（$Y^{a=0}=Y^{a=1}=0$）。由于随机抽样与随机化，以下三组具有相同的反事实类型分布：试验中受治疗的 1000 人、试验中未受治疗的 1000 人，以及观察性研究中的 8000 人。关键的观察是：观察性研究中 0\% 的死亡率意味着 (i) 不存在“注定者”；(ii) 所有 $Y^{a=1}=1$ 的个体必属“受害”型且接受了 $A=0$；(iii) 所有 $Y^{a=0}=1$ 的个体必属“受益”型且接受了 $A=1$。

接下来我们用这些观察，按反事实类型重建试验数据。如上所论，550 名死亡（$Y^{a=1}=1$）的 $A=1$ 个体属于“受害”型。由随机化，必定还有 550 名 $A=0$ 的“受害”型个体，他们当然存活了。同理，我们可以填出“受益”型的表。

\begin{center}
\begin{tabular}[t]{lcc}
\toprule
\multicolumn{3}{c}{“受益”型}\\
& $A=0$ & $A=1$\\
\midrule
$Y=1$ & 450 & 0\\
$Y=0$ & 0 & 450\\
\midrule
& 450 & 450\\
\bottomrule
\end{tabular}
\hspace{3em}
\begin{tabular}[t]{lcc}
\toprule
\multicolumn{3}{c}{“受害”型}\\
& $A=0$ & $A=1$\\
\midrule
$Y=1$ & 0 & 550\\
$Y=0$ & 550 & 0\\
\midrule
& 550 & 550\\
\bottomrule
\end{tabular}
\end{center}

把两张表合并，我们就恢复了试验的总体数据，这意味着不存在“免疫”型个体。我们的结论是：监管机构应当重新许可处理 $A$，并建议维持现行做法，因为观察性研究显示，每个肺癌患者都以某种方式私下知道——这是试验者无从获得的——该处理对自己是有害还是有益。下面我们描述一个极端情形，它可以解释这种私下知识的来源。

假设人群中存在 55/45 的族群划分，且由于遗传原因，该处理对第一族群中的肺癌患者一致有害（即全属“受害”型），而对第二族群中的所有癌症患者一致有益（即全属“受益”型）。再假设在试验之前的某个时候，第一族群中的癌症患者在看到本族群若干成员服药后迅速死亡之后，便不再服用该处理。相反，假设第二族群中的所有癌症患者在看到本族群若干成员服药后存活下来之后，都选择服用该处理。换言之，假设观察性数据中同时存在 (i) 由族群导致的极大定性效应修饰，与 (ii) 由族群导致的极大混杂。

上述极端设定当然不现实，但它有助于说明一个重要观点：随机化试验比较的是“治疗所有人”与“不治疗任何人”这两种策略，而最优策略是“只治疗第二族群中的个体”。当未获取效应修饰因子（即族群）的数据时，随机化试验无法把个体分配到最优策略。相反，在观察性研究中，所有人都遵循了最优策略，因而得到了无人死亡的最优结局。于是，揭示了正确政策的，是有混杂的观察性研究，而不是无混杂的随机化试验。
\end{finepoint}

\begin{finepoint}[把精要 22.7 推广到现实设定]
上述讨论可以推广到更现实的设定，以说明：把随机化试验数据与观察性数据结合起来的设计，可能比仅有随机化试验数据的设计更有信息量——前提是随机化数据与观察性数据中的个体都是所有符合试验条件者的随机样本。

假设我们在一项二值处理 $A$ 与结局 $Y$ 上开展随机化试验，而该人群中处理已在使用（$Y$ 取值越低越好）。再假设，如偶尔会发生的那样，观察性研究中的平均结局 $\E[Y]$ 低于随机化试验两个臂的平均结局，这意味着 $\E[Y]$ 同时低于 $\E[Y^{a=0}]$ 与 $\E[Y^{a=1}]$。于是，我们也许会选择让当前社区的处理实践保持不变。

下面我们证明：对某个不同于试验中所比较的“治疗所有人”（$a=1$）与“不治疗任何人”（$a=0$）的策略 $g$，有 $\E[Y]=\E[Y^{g}]$。具体而言，令 $U$ 为一组充分的、未测量的（可能未知的）处理前协变量，足以确保 $Y^{a}\indep A\given U$；并令 $\Pr[A=1\given U]$ 为观察性数据中相应的倾向性评分。那么，对处理 $A=1$ 以概率 $\Pr[A=1\given U]$ 随机分配的策略 $g$，上述等式成立，因为这一选择恰好给出了生成 $A$、$Y$ 与未知 $U$ 之观察性数据的随机策略 $g$。尽管由于 $U$ 的数据不可得而无法实施策略 $g$，上述讨论仍可促使研究者测量处理前协变量 $V$，用以分析随机化试验数据，从而找出并实施一个确定性动态策略 $g^{*}(x)$，使由无混杂的随机化试验数据估计的 $\E[Y^{g^{*}}]$ 低于观察性数据中的 $\E[Y]$。
\end{finepoint}
